%=============================================================================
% UIDT v3.9 — EFFECTIVE GAP DERIVATION DOSSIER
% An Effective Phenomenological Derivation of the Yang-Mills Scale
% MODULAR VERSION - Uses \input for Appendices
%=============================================================================
% Author: Philipp Rietz (ORCID: 0009-0007-4307-1609)
% DOI: 10.5281/zenodo.17835200
% License: CC BY 4.0
% Version: 3.9 (May 2026)
%=============================================================================
% 
% NOTE: This is the MODULAR version. 
% For the fully integrated version, use main-complete.tex
%
% Structure:
%   main.tex (this file) - Preamble, Main Text, Bibliography
%   UIDT_Appendix_A_OS_Axioms.tex - OS Axioms detailed proofs
%   UIDT_Appendix_B_BRST.tex - BRST Cohomology
%   UIDT_Appendix_C_Numerical.tex - Numerical Verification
%   UIDT_Appendix_D_Auxiliary.tex - Auxiliary Field Elimination
%   UIDT_Appendix_G_Extended.tex - Extended Proofs (GNS, Spectral, etc.)
%   UIDT_Appendix_H_GapAnalysis.tex - Clay Gap Analysis Summary
%   UIDT_Appendix_I_LightQuarkMasses_Torsion.tex - Light Quark Torsion
%   UIDT_Chapter_7_QuarkMassHierarchy.tex - Quark Mass Hierarchy
%   UIDT_v3.9_QuarkMass_Supplement.tex - Quark Mass Supplement
%=============================================================================

\documentclass[11pt,a4paper]{article}

%-----------------------------------------------------------------------------
% PACKAGES
%-----------------------------------------------------------------------------
\usepackage[utf8]{inputenc}
\usepackage[T1]{fontenc}
\usepackage{lmodern}
\usepackage{amsmath,amssymb,amsthm}
\usepackage{mathtools}
\usepackage{physics}
\usepackage{bm}
\usepackage{bbold}
\usepackage{graphicx}
\usepackage{xcolor}
\usepackage{hyperref}
\usepackage{cleveref}
\usepackage{booktabs}
\usepackage{longtable}
\usepackage{enumitem}
\usepackage{tcolorbox}
\usepackage{fancyhdr}
\usepackage{geometry}
\usepackage{titlesec}
\usepackage{appendix}

\geometry{margin=2.5cm}

%-----------------------------------------------------------------------------
% CUSTOM COMMANDS
%-----------------------------------------------------------------------------
\newcommand{\R}{\mathbb{R}}
\newcommand{\C}{\mathbb{C}}
\newcommand{\N}{\mathbb{N}}
\newcommand{\Z}{\mathbb{Z}}
\newcommand{\Hilbert}{\mathcal{H}}
\newcommand{\Schwinger}{\mathcal{S}}
\newcommand{\Lagr}{\mathcal{L}}
\newcommand{\im}{\mathrm{im}\,}
\newcommand{\ke}{\mathrm{ker}\,}
\newcommand{\Tr}{\mathrm{Tr}}
\newcommand{\supp}{\mathrm{supp}}
\newcommand{\spec}{\mathrm{spec}}

%-----------------------------------------------------------------------------
% THEOREM ENVIRONMENTS
%-----------------------------------------------------------------------------
\theoremstyle{plain}
\newtheorem{theorem}{Theorem}[section]
\newtheorem{lemma}[theorem]{Lemma}
\newtheorem{proposition}[theorem]{Proposition}
\newtheorem{corollary}[theorem]{Corollary}

\theoremstyle{definition}
\newtheorem{definition}[theorem]{Definition}
\newtheorem{axiom}[theorem]{Axiom}
\newtheorem{remark}[theorem]{Remark}
\newtheorem{example}[theorem]{Example}

%-----------------------------------------------------------------------------
% PAGE FORMATTING
%-----------------------------------------------------------------------------
\clubpenalty=10000
\widowpenalty=10000
\displaywidowpenalty=10000
\raggedbottom

%-----------------------------------------------------------------------------
% HEADERS
%-----------------------------------------------------------------------------
\pagestyle{fancy}
\fancyhf{}
\fancyhead[L]{\small UIDT v3.9 --- Effective Yang-Mills Scale}
\fancyhead[R]{\small Phenomenological Derivation Dossier}
\fancyfoot[C]{\thepage}

%-----------------------------------------------------------------------------
% TITLE
%-----------------------------------------------------------------------------
\title{%
\vspace{-2cm}
{\large UIDT Framework v3.9 --- Effective Yang-Mills Scale Dossier}\\[0.5cm]
\textbf{An Effective Phenomenological Derivation of the Yang-Mills Scale}\\[0.3cm]
{\large via Hybrid Information-Density Projection}\\[0.5cm]
{\normalsize Evidence-tagged version; no Clay-level proof claim}
}

\author{%
Philipp Rietz\\
\small ORCID: 0009-0007-4307-1609\\
\small DOI: 10.5281/zenodo.17835200
}

\date{May 2026}

%=============================================================================
\begin{document}
%=============================================================================

\maketitle
\thispagestyle{empty}

\begin{abstract}
We present an effective phenomenological derivation of a Yang-Mills spectral
scale in the UIDT v3.9 framework. The construction introduces an effective
information-density scalar field $S(x)$ and treats the infrared sector through
a mean-field projection. Within this reduced algebraic model, the gap equation
defines a local contraction with Lipschitz constant
$L=3.749\times10^{-5}\ll1$ and a unique fixed point
$\Delta^*=1.710\pm0.015\,\mathrm{GeV}$ [A, internal reduced-model closure].
The mapping from this one-dimensional contraction to the full
infinite-dimensional pure Yang-Mills path integral is not proven and remains
an explicit Stratum III assumption [E]. The calibrated kinetic parameter
$\gamma=16.339$ remains [A-], not RG-derived. Accordingly, this dossier does
not claim a Clay Mathematics Institute proof of the Yang-Mills existence and
mass gap problem. Its purpose is to document the hybrid method, numerical
closure, lattice-compatible scale comparison, and falsifiable follow-up
predictions without evidence inflation.
\end{abstract}

\tableofcontents
\newpage

%=============================================================================
% PART I: MAIN TEXT
%=============================================================================

\part{Effective Gap Derivation}

%-----------------------------------------------------------------------------
\section{Introduction: Yang-Mills Scale and Epistemic Boundary}
\label{sec:introduction}
%-----------------------------------------------------------------------------

\subsection{Statement of the Problem}

The Clay Mathematics Institute Millennium Prize Problem concerning 
Yang-Mills theory requires:
\begin{quote}
\textit{Prove that for any compact simple gauge group $G$, a non-trivial 
quantum Yang-Mills theory exists on $\R^4$ and has a mass gap $\Delta > 0$.}
\end{quote}

This dossier does not claim to satisfy that standard. It documents an
effective UIDT construction whose reduced algebraic sector is internally
closed, while the pure-theory equivalence remains open.

\subsection{Structure of the Hybrid Derivation}

The hybrid derivation proceeds through the following steps:
\begin{enumerate}
\item Definition of the augmented Lagrangian with scalar field $S(x)$
\item Effective infrared mean-field projection [E]
\item Reduced algebraic gap equation [A within the model]
\item Local Banach fixed-point closure [A within the model]
\item Calibration and lattice-compatible scale comparison [A-]/[B]
\item Explicit documentation of the pure Yang-Mills boundary [E]
\item Falsifiable tensor-glueball and finite-temperature predictions [D]
\end{enumerate}

%-----------------------------------------------------------------------------
\section{Axiomatic Framework: Fields and Lagrangian}
\label{sec:fields}
%-----------------------------------------------------------------------------

\subsection{Field Content}

The effective theory contains:
\begin{itemize}
\item $A^a_\mu(x)$: $\mathrm{SU}(3)$ gauge field ($a = 1,\ldots,8$)
\item $S(x)$: Real scalar field, treated here as an effective information-density field
\item $c^a(x), \bar{c}^a(x)$: Faddeev-Popov ghosts
\item $B^a(x)$: Nakanishi-Lautrup auxiliary field
\end{itemize}

\subsection{The Lagrangian Density}

\begin{definition}[UIDT Effective Lagrangian]
\label{def:lagrangian}
The complete gauge-fixed effective Lagrangian is:
\begin{equation}
\Lagr = \Lagr_{\mathrm{YM}} + \Lagr_S + \Lagr_{\mathrm{coupling}} 
+ \Lagr_{\mathrm{gf}} + \Lagr_{\mathrm{ghost}}
\label{eq:full_lagrangian}
\end{equation}
with components:
\begin{align}
\Lagr_{\mathrm{YM}} &= -\frac{1}{4}F^a_{\mu\nu}F^{a\mu\nu} \\
\Lagr_S &= \frac{1}{2}(\partial_\mu S)^2 - V(S) \\
\Lagr_{\mathrm{coupling}} &= \frac{\kappa}{\Lambda}S\,\Tr(F_{\mu\nu}F^{\mu\nu}) \\
\Lagr_{\mathrm{gf}} &= B^a(\partial_\mu A^{a\mu}) + \frac{\xi}{2}(B^a)^2 \\
\Lagr_{\mathrm{ghost}} &= \bar{c}^a \partial_\mu D_\mu^{ab} c^b
\end{align}
\end{definition}

The scalar potential is:
\begin{equation}
V(S) = \frac{1}{2}m_S^2 S^2 + \frac{\lambda_S}{4!}S^4
\label{eq:potential}
\end{equation}

%-----------------------------------------------------------------------------
\section{Osterwalder-Schrader Axioms}
\label{sec:os}
%-----------------------------------------------------------------------------

\begin{theorem}[Reduced-Model OS Consistency Check]
The Schwinger functions of the UIDT effective model are required to satisfy the
Osterwalder-Schrader consistency conditions: Temperedness (OS0), Euclidean
Covariance (OS1), Symmetry (OS2), Cluster Property (OS3), and Reflection
Positivity (OS4). This statement is part of the reduced-model consistency
program and is not a substitute for a completed construction of the pure
Yang-Mills measure.
\end{theorem}

[Detailed checks in Appendix A]

%-----------------------------------------------------------------------------
\section{Wightman Reconstruction}
\label{sec:wightman}
%-----------------------------------------------------------------------------

\begin{theorem}[Reduced-Model OS-Wightman Reconstruction]
If the reduced effective model satisfies the OS axioms, Wightman reconstruction
applies to that effective model. This does not by itself establish the same
result for pure Yang-Mills theory.
\end{theorem}

%-----------------------------------------------------------------------------
\section{BRST Cohomology}
\label{sec:brst}
%-----------------------------------------------------------------------------

\begin{theorem}[BRST Nilpotency]
\label{thm:nilpotency}
The BRST operator $s$ satisfies $s^2 = 0$ within the effective gauge-fixed
model.
\end{theorem}

\begin{definition}[Physical Hilbert Space]
\begin{equation}
\Hilbert_{\mathrm{phys}} = \ke Q / \im Q
\label{eq:physical_hilbert}
\end{equation}
\end{definition}

[Detailed treatment in Appendix B]

%-----------------------------------------------------------------------------
\section{Gauge Independence}
\label{sec:gauge}
%-----------------------------------------------------------------------------

\begin{theorem}[Nielsen Identity]
\label{thm:nielsen}
Physical observables of the effective model must be independent of the gauge parameter $\xi$:
\begin{equation}
\frac{\partial\Delta^*}{\partial\xi} = 0
\end{equation}
\end{theorem}

%-----------------------------------------------------------------------------
\section{Renormalization Group Invariance}
\label{sec:rg}
%-----------------------------------------------------------------------------

\begin{theorem}[UV Fixed Point Constraint]
\label{thm:fixed_point}
The reduced UIDT closure uses a fixed-point constraint satisfying:
\begin{equation}
5\kappa^{*2} = 3\lambda_S^*
\end{equation}
with $\kappa^* = 0.500 \pm 0.008$.
\end{theorem}

%-----------------------------------------------------------------------------
\section{Effective Phenomenological Derivation of the Yang-Mills Scale}
\label{sec:massgap}
%-----------------------------------------------------------------------------

\subsection{The Effective Mean-Field Projection}

Instead of claiming a global proof over the full infinite-dimensional
functional space of pure Yang-Mills theory, we formulate an effective infrared
projection. We assume that the non-perturbative vacuum dynamics can be
parameterized by coupling to an effective scalar information-density field
$S(x)$.

Under the additional assumption that the infrared system is dominated by the
model condensate input $\mathcal C$, the path integral is reduced to an
algebraic gap equation. The exact topological equivalence between this
reduction and pure Yang-Mills theory remains an open conjectural step [E].

\subsection{The Algebraic Gap Equation and Local Banach Iteration}

Within the reduced effective model, define
\[
T(\Delta)=
\sqrt{
m_S^2+
\frac{\kappa^2\mathcal C}{4\Lambda^2}
\left[
1+\frac{\ln(\Lambda^2/\Delta^2)}{16\pi^2}
\right]
}.
\]

The deterministic fixed-point constraint is
\[
5\kappa^2=3\lambda_S,
\qquad
\lambda_S=\frac{5}{12}
\quad\text{for}\quad
\kappa=\frac12.
\]
The residual is therefore exactly zero in rational arithmetic and must remain
below $10^{-14}$ in numerical verification.

The local derivative of the reduced operator is
\[
T'(\Delta)
=-\frac{\alpha\beta}{\Delta T(\Delta)},
\qquad
\alpha=\frac{\kappa^2\mathcal C}{4\Lambda^2},
\quad
\beta=\frac{1}{16\pi^2}.
\]
Near the fixed point,
\[
L\simeq 3.749\times10^{-5}\ll1.
\]

Consequently, $T$ is a strong local contraction for the reduced one-dimensional
model interval. Banach iteration gives
\[
\Delta^*
=1.710035046742213182020771096614\ldots\,\mathrm{GeV}.
\]
This is an [A] statement about the reduced model only.

\subsection{Phenomenological Calibration}

The numerical scale is not a purely ab-initio prediction from the
dimensionless Yang-Mills coupling $g$. It depends on the external/model
condensate input $\mathcal C$ and on the calibrated kinetic parameter
$\gamma=16.339$ [A-]. The closure demonstrates internal consistency of the
hybrid parameter system, not a complete proof of the Clay problem.

\subsection{Acknowledged Limitations}

\begin{enumerate}
\item The contraction proof is local and one-dimensional; it is not a proof on
      the full measure of the infinite-dimensional QFT path integral.
\item Removing the auxiliary/effective field $S(x)$ by a smooth deformation is
      not trivial in QFT. Haag-type obstructions and phase transitions cannot
      be excluded by the current argument.
\item The use of condensate inputs classifies the construction as an effective
      phenomenological model, not as a pure axiomatic derivation from $g$ alone.
\end{enumerate}

%-----------------------------------------------------------------------------
\section{Auxiliary Field Boundary}
\label{sec:auxiliary}
%-----------------------------------------------------------------------------

\begin{remark}[Auxiliary Field Boundary]
The scalar field $S(x)$ may be integrated out in the reduced effective model to
generate induced Yang-Mills-sector operators. This does not prove that the
resulting effective action is continuously equivalent to pure Yang-Mills theory
without phase transition or spectral discontinuity.
\end{remark}

[Details in Appendix D]

%-----------------------------------------------------------------------------
\section{Lattice QCD Comparison}
\label{sec:lattice}
%-----------------------------------------------------------------------------

\begin{theorem}[Lattice Compatibility]
\label{thm:compatibility}
The UIDT reduced-model scale $\Delta^* = 1.710 \pm 0.015\,\mathrm{GeV}$ is
compatible with selected lattice QCD determinations under the stated
normalization and uncertainty conventions. This is a [B] comparison statement,
not a proof of the pure Yang-Mills spectrum.
\end{theorem}

%-----------------------------------------------------------------------------
\section{Conclusion}
\label{sec:conclusion}
%-----------------------------------------------------------------------------

We have documented an effective phenomenological derivation of a Yang-Mills
spectral scale within UIDT v3.9. The reduced gap equation is internally closed
and numerically stable, but this does not constitute a Clay-level proof of pure
Yang-Mills existence and mass gap.

\subsection{Methodological Limitations}

\begin{enumerate}
\item \textbf{Scalar Field Extension:} The construction employs an effective
      scalar field $S(x)$ not present in pure Yang-Mills. The connection to pure
      Yang-Mills remains an open projection-equivalence problem [E].
\item \textbf{Gauge Group:} The current numerical scale is stated for the
      $\mathrm{SU}(3)$ sector.
\item \textbf{Deformation Limit:} The continuous deformation to pure Yang-Mills
      is not formalized as rigorous homotopy and cannot be used as a Clay-level
      conclusion.
\end{enumerate}

\input{UIDT_Chapter_7_QuarkMassHierarchy}

%=============================================================================
% PART II: APPENDICES
%=============================================================================
\newpage
\part{Mathematical Appendices}

\appendix

% Include separate appendix files
\input{UIDT_Appendix_A_OS_Axioms}
\input{UIDT_Appendix_B_BRST}
\input{UIDT_Appendix_C_Numerical}
\input{UIDT_Appendix_D_Auxiliary}
\input{UIDT_Appendix_G_Extended}
\input{UIDT_Appendix_H_GapAnalysis}
\input{UIDT_Appendix_I_LightQuarkMasses_Torsion}
\input{UIDT_v3.9_QuarkMass_Supplement}
%=============================================================================
% BIBLIOGRAPHY
%=============================================================================
\newpage
\begin{thebibliography}{99}

\bibitem{JaffeWitten}
A. Jaffe and E. Witten, ``Quantum Yang-Mills Theory,'' 
Clay Mathematics Institute, 2000.

\bibitem{Morningstar1999}
C. Morningstar and M. Peardon, 
``The glueball spectrum from an anisotropic lattice study,''
Phys. Rev. D \textbf{60}, 034509 (1999).

\bibitem{Chen2006}
Y. Chen \textit{et al.}, 
``Glueball spectrum and matrix elements,''
Phys. Rev. D \textbf{73}, 014516 (2006).

\bibitem{OsterwalderSchrader}
K. Osterwalder and R. Schrader, 
``Axioms for Euclidean Green's functions,''
Commun. Math. Phys. \textbf{31}, 83 (1973).

\bibitem{KugoOjima}
T. Kugo and I. Ojima, 
``Local covariant operator formalism,''
Prog. Theor. Phys. Suppl. \textbf{66}, 1 (1979).

\bibitem{Nielsen1975}
N. K. Nielsen, 
``On the gauge dependence of spontaneous symmetry breaking,''
Nucl. Phys. B \textbf{101}, 173 (1975).

\bibitem{Wetterich1993}
C. Wetterich, 
``Exact evolution equation for the effective potential,''
Phys. Lett. B \textbf{301}, 90 (1993).

\bibitem{SVZ}
M. A. Shifman, A. I. Vainshtein, and V. I. Zakharov,
``QCD and resonance physics,''
Nucl. Phys. B \textbf{147}, 385 (1979).

\bibitem{BRSTBecchi}
C. Becchi, A. Rouet, and R. Stora,
``Renormalization of gauge theories,''
Ann. Phys. \textbf{98}, 287 (1976).

\bibitem{Tyutin}
I. V. Tyutin,
``Gauge invariance in field theory,''
Lebedev preprint FIAN-39 (1975).

\end{thebibliography}

\end{document}
